\section{Discussion}

\subsection{Implications for UK–Sri Lanka Transnational Delivery}

For UK providers delivering programmes with Sri Lankan partners, alignment with the 
Frameworks for Higher Education Qualifications (FHEQ) and, where relevant, SCQF equivalences 
is essential for maintaining academic coherence, transparent progression, and credible 
recognition. Programme structures should map explicitly to level descriptors, learning 
outcomes, and typical credit volumes, ensuring that students can understand how their 
programme fits within the wider UK qualifications system.

Under the Office for Students (OfS) regulatory framework, the registered lead provider retains 
full responsibility for academic standards and quality wherever provision is delivered. This 
includes validated, franchised, subcontracted, or distance-learning delivery in Sri Lanka. 
Therefore, local teaching and assessment must operate within the awarding provider’s academic 
regulations, quality cycle, and reporting structures. Student protection planning, consumer 
law compliance, data duties, and reportable-events handling remain obligations of the UK 
provider; Sri Lankan delivery partners must support these requirements operationally.

Recognition and admissions processes should draw on UK~ENIC Statements of Comparability to 
ensure defensible decisions for both entry and advanced standing. This underpins transparency, 
equity, and consistency for Sri Lankan applicants, particularly where qualification mapping 
between SLQF and FHEQ levels is required. Sector evidence from British Council Sri Lanka and 
Universities UK International highlights areas of opportunity and ongoing risk signals in local 
TNE delivery, informing due diligence, governance, and student-outcome monitoring frameworks.

\begin{tcolorbox}
\textbf{Key implications for UK--Sri Lanka provision:}
\begin{itemize}
    \item Programme design should follow FHEQ level descriptors and UK credit norms.
    \item Exit-award pathways must be transparent and mapped to FHEQ and SLQF structures.
    \item The UK lead provider remains wholly accountable under OfS Conditions A--G.
    \item Admissions and advanced standing require UK~ENIC-supported comparability.
    \item Local partners must support student protection, data, consumer, and reporting duties.
    \item TNE risk signals from British Council Sri Lanka and UUK should shape partnership oversight.
\end{itemize}
\end{tcolorbox}

\subsection{Limitations and Practical Constraints}

Interpretation of regulatory expectations across the UK home nations is influenced by 
structural differences. Scotland’s enhancement-led quality model, governed through SFC 
guidance, is not directly equivalent to the OfS registration system. Wales, through the 
Commission for Tertiary Education and Research (CTER), integrates oversight into a single 
tertiary framework, including policy for registration, designation, and quality assessment. 
Northern Ireland’s Department for the Economy provides policy and funding oversight rather than 
a codified provider-registration framework. As a result, comparisons must be made carefully, 
acknowledging the different levers and regulatory instruments in each nation.

Further limitations include website restructuring between evidence capture and analysis, 
availability of PDFs without stable pagination, and the descriptive rather than prescriptive 
nature of many sector reports. These affect generalisability and require professional 
interpretation, especially when translating regulatory models to transnational delivery 
contexts.

\begin{tcolorbox}
\textbf{Key limitations:}
\begin{itemize}
    \item Variation in regulatory structures across UK nations affects comparability.
    \item SFC, CTER, and DfE~NI issue guidance that differs from OfS conditions in scope.
    \item Some evidence (e.g., TNE reports) is descriptive, not regulatory.
    \item Website changes and document availability may affect reproducibility.
    \item Local SLQF alignment requires careful judgement when interpreting UK standards.
\end{itemize}
\end{tcolorbox}
